\section{Выводы}

В ходе выполнения лабораторных работ была реализована структура данных PATRICIA
для словаря строковых ключей и проведено исследование качества полученной программы.
Реализованы поиск, вставка, удаление, сохранение и загрузка словаря из бинарного файла.

Профилирование с помощью \texttt{gprof} показало, что основное время расходуется в
операциях \texttt{Remove}, \texttt{Find} и \texttt{Insert}, что соответствует устройству
программы: именно эти функции выполняют обход дерева по критическим битам.

Проверка Valgrind Memcheck не обнаружила утечек памяти и некорректных обращений в коде
словаря. Massif показал умеренное потребление памяти: около 7.335 МБ heap-памяти для
набора из $100\,000$ слов.

Были исправлены функциональные недочёты сериализации: загрузка отсутствующего файла,
повреждённого файла и ключа максимальной длины 256 символов. Также был исправлен benchmark,
чтобы измерение поиска PATRICIA не могло быть сведено оптимизатором к пустому циклу.

Главный результат работы~--- программа стала надёжнее на ошибочных входах и честнее
измеряется в тестах производительности, при этом базовая структура данных и её
асимптотические свойства остались прежними.

\pagebreak
